\section{Experimental Setup}
\label{sec:exp}

\textbf{Data and preprocessing.} We use the dataset of Section~\ref{sec:data} (4{,}563 main-tank samples
from 628 transformers, 2019--2024; median 6 samples per unit over $\sim$4.5 years). Gas values
are clipped at the 99.9th percentile, missing readings imputed as zero (below detection),
log-transformed, and standardised; SD-CAE additionally uses the CLR composition of the five
combustible gases (H2, CH4, C2H2, C2H4, C2H6). The weak label \texttt{fault\_note} is positive for 71 of 628 units
($11.3\%$).

\textbf{Baselines.} For fault-type recovery we compare the deployed CLR-PCA representation
against clustering on standardised log-concentrations (the severity baseline) and on raw
proportions, KMeans and a Gaussian mixture directly on the 5-D CLR, and the SD-CAE autoencoder
with adversary weights $\lambda\in\{0,1,10\}$ (all negative ablations). For risk screening we compare the
index against conventional severity proxies---total combustible gas, a C57.104-style count of
gases above the fleet 90th percentile, and a Duval-arcing flag---and against the sample-count
confound. For the anomaly term we compare autoencoder reconstruction, Isolation Forest, and
Deep~SVDD, plus a battery of standard detectors.

\textbf{Metrics.} Fault-type recovery is measured by ARI and NMI against the Duval classes (a
deterministic rule, hence a confound-free target). Risk ranking is measured by AUC against
\texttt{fault\_note}, precision at the top-$k$ units, and the event rate by risk decile;
severity leakage of a representation is measured by $R^2(m\mid\mathbf{z})$. Stochastic results
are reported as mean~$\pm$~standard deviation over multiple seeds; headline AUCs carry
unit-level bootstrap $95\%$ confidence intervals, paired comparisons use a paired bootstrap,
and the chemistry target of Section~\ref{sec:results} uses a unit-blocked cluster bootstrap
because its positives concentrate in few units.

\textbf{Implementation.} Models are small multilayer perceptrons trained on CPU (PyTorch); the
data is small enough that no GPU is required. Every figure and number is produced by versioned
scripts writing to a results directory, with fixed seeds for reproducibility. The main
hyperparameters are listed in Table~\ref{tab:hp}.

\begin{table}[t]
\centering
\caption{Main hyperparameters.}
\label{tab:hp}
\begin{tabular}{ll}
\toprule
Component & Setting \\
\midrule
Deployed representation & CLR $\rightarrow$ standardise $\rightarrow$ PCA, 2 components \\
SD-CAE (ablation) & MLP $16\text{--}8$, ReLU, latent 2; Adam, lr $10^{-3}$ \\
SD-CAE training & batch 128, 200 epochs, early stop (patience 20) \\
CLR zero-replacement $\delta$ & $10^{-3}$ (sensitivity: $10^{-4}$--$10^{-2}$) \\
Adversary weight $\lambda$ & ablated at 0, 1, 10 (not deployed) \\
Index weights & severity 1, acetylene 2, temporal 1, anomaly 1 \\
Clustering & KMeans, $k=7$ (Duval classes); silhouette $k$ as check \\
\bottomrule
\end{tabular}
\end{table}
